Die Frage nach den Zahlkörpern kleinsten Grades über einem Grundkörper $P$, in denen eine in $P$ irreduzible Darstellung einer endlichen Gruppe in absolut irreduzible Bestandteile zerfällt, wurde zuerst von I. Schur behandelt.\footnote{Arithmetische Untersuchungen über endliche Gruppen, Sitzungsber. d. Berl. Ak. d. Wiss. 1906, S. 164. Beiträge zur Theorie der Gruppen linearer Substitutionen, Transact. of the Am. Math. Soc. Ser. 2, Vol. XV, 1909, S. 159. In der zweiten Arbeit werden die Ergebnisse auf vollständig reduzible hyperkomplexe Systeme übertragen. In eine Klasse wird eine Darstellung mit all ihren Transformierten (ähnlichen) zusammengefaßt.} Es enthalte $P$ etwa den Charakter eines derartigen Bestandteiles. Dann zeigte I. Schur, daß der Grad dieser Körper in bezug auf $P$ -- der Index -- übereinstimmt mit der Anzahl jener absolut irreduziblen Bestandteile, die alle derselben Klasse angehören; daß folglich derselbe Körper schon bestimmt ist durch die Forderung, einen absolut irreduziblen Bestandteil abzutrennen; ferner daß der Grad jedes Körpers in bezug auf $P$, in dem eine solche Zerfällung statthat, ein Vielfaches des Index wird. Alle diese Körper sollen als Zerfällungskörper bezeichnet werden, auch für den Fall, daß $P$ den Charakter nicht enthält. An einem Beispiel zeigte I. Schur, daß die Zerfällungskörper kleinsten Grades nicht isomorph zu sein brauchen. Enthielt $P$ keinen zugehörigen einfachen Charakter, so müssen die Zerfällungskörper den Körper eines solchen und seine Konjugierten in bezug auf $P$ umfassen; die zu verschiedenen konjugierten Charakteren gehörigen absolut irreduziblen Darstellungen sind zwar konjugiert, gehören aber offenbar zu verschiedenen Klassen. Zur Abtrennung eines Systems absolut irreduzibler Darstellungen einer Klasse genügt auch hier der Körper des zugehörigen Charakters und einer der entsprechenden Zerfällungskörper.
